\chapter{Orthogonal Range Searching}
\label{ch:range_searching}

\section{Range Queries}
\label{sec:rangesearch_queries}

A \textbf{range query}\index{range query} asks for the points of a fixed set
$P \subset \mathbb{R}^d$ that lie inside a query region $Q$: the output is
either the \emph{count} $|P \cap Q|$ or the \emph{report} of all points in
$P \cap Q$. The queries are answered by a data structure built once from
$P$, and the structure is judged by its storage, its query time, and its
construction and update costs. The problem is fundamental because it is the
primitive behind database range selection, geographic windowing,
multi-dimensional indexing, and the nearest-neighbor queries of
Chapter~\ref{ch:proximity}.

The query region $Q$ determines the difficulty of the problem:
\begin{description}
  \item[Orthogonal ranges] $Q$ is an axis-aligned box (a rectangle in the
    plane, a box in $\mathbb{R}^d$), possibly with only some sides
    specified (a 3-sided rectangle, a half-plane strip). These queries
    support the elegant tree structures of this chapter.
  \item[Simplex ranges] $Q$ is a half-space or a simplex; these queries
    require the partition-tree and cutting techniques of
    Section~\ref{sec:rangesearch_simplex}.
  \item[General ranges] $Q$ is a disk, a polygon, or a set of several
    ranges; these are typically answered by reduction to the orthogonal or
    simplex case.
\end{description}

This chapter develops the classical orthogonal range-searching hierarchy:
one-dimensional search (Section~\ref{sec:rangesearch_1d}), $k$-d trees
(Section~\ref{sec:rangesearch_kd_trees}), range trees with multilevel
structures (Sections~\ref{sec:rangesearch_range_trees}--\ref{sec:rangesearch_multilevel}),
fractional cascading (Section~\ref{sec:rangesearch_fractional_cascading}),
priority search trees (Section~\ref{sec:rangesearch_priority_search}), and
the higher-dimensional and dynamic extensions
(Sections~\ref{sec:rangesearch_higher_dimensional}--\ref{sec:rangesearch_dynamic}).
The statistical structure of points in high dimensions is treated in
Section~\ref{sec:rangesearch_curse} and
Chapter~\ref{ch:high_dimensional}.

\section{One-Dimensional Range Searching}
\label{sec:rangesearch_1d}

\subsection{Overview}

In one dimension the points are a set of $n$ real numbers and the query is
an interval $[x_1, x_2]$: report (or count) the points in the interval. This
is the simplest nontrivial range-searching problem, and the ideas it
introduces --- sorted order, binary search, and split nodes --- are the
basis of every structure in this chapter.

\subsection{Definitions}

\begin{definition}[Range Counting and Reporting]\label{def:rangesearch_counting}
  For a set $P$ of $n$ reals and an interval $I=[x_1,x_2]$, the
  \textbf{counting query} asks for $|P \cap I|$ and the \textbf{reporting
  query} asks for the sorted list of points in $P \cap I$.
\end{definition}

\subsection{Theory}

Sort the points as $p_1 < p_2 < \cdots < p_n$. The points in the query
interval form a contiguous block $p_{i+1}, \dots, p_j$, where $i$ is the
number of points $\le x_1$ and $j$ the number of points $< x_2$. Both
indices are found by binary search on the sorted array. A balanced binary
search tree keyed by value stores the same information with the additional
ability to support insertion and deletion (Section~\ref{sec:rangesearch_dynamic}).

\subsection{Pseudocode}

\begin{algorithm}[htbp]
\caption{One-Dimensional Range Query}
\label{alg:rangesearch_1d}
\begin{algorithmic}[1]
\Require Sorted array $A[1..n]$ of distinct points, query interval $[x_1,x_2]$.
\Ensure The points of $A$ in $[x_1,x_2]$.
\Function{RangeQuery}{$A, x_1, x_2$}
  \State $i \gets \text{binarySearch}(A, x_1)$  \Comment{first index with $A[i] \ge x_1$}
  \State $j \gets \text{binarySearch}(A, x_2)$  \Comment{last index with $A[j] \le x_2$}
  \State \Return $A[i..j]$
\EndFunction
\end{algorithmic}
\end{algorithm}

\subsection{Complexity}

\begin{itemize}
  \item \textbf{Time}: counting in $O(\log n)$; reporting in
    $O(\log n + k)$ with $k = |P \cap I|$ (the output itself needs $k$ time).
  \item \textbf{Storage}: $O(n)$.
\end{itemize}

\subsection{Usages}

\begin{itemize}
  \item \textbf{Database queries}: selecting records whose key lies in an
    interval.
  \item \textbf{Component of higher-dimensional structures}: every range
    tree node stores a one-dimensional structure of exactly this kind
    (Section~\ref{sec:rangesearch_range_trees}).
\end{itemize}

\subsection{References}

The one-dimensional structure is elementary; its textbook treatment
appears in \cite{berg2008} and the analysis of its use inside multilevel
structures in \cite{bentley1980multidim}.

\section{k-d Trees}
\label{sec:rangesearch_kd_trees}

\subsection{Overview}

The \textbf{$k$-d tree}\index{$k$-d tree} partitions the point set
recursively with axis-aligned splitting hyperplanes: at each internal node,
the points are split by a hyperplane perpendicular to one coordinate axis
into two subsets of equal cardinality, each stored in a child. The result
is a balanced binary tree whose leaves are points. $k$-d trees support
orthogonal range queries by pruning any subtree whose cell is disjoint from
the query range. A brief outline appears here; the full data structure ---
construction, balancing, and query pruning --- is developed in
Section~\ref{sec:kd_tree} (Chapter~\ref{ch:spatial_structures}).

\subsection{Definitions}

\begin{definition}[Split Plane and Cell]\label{def:kd_cell}
  Each $k$-d tree node $v$ stores a \textbf{split hyperplane} $x_d = c$
  (perpendicular to the split axis $d$) and is associated with a
  \textbf{cell}, the axis-aligned box containing the points of its subtree.
  The two children own the two halves of the cell on either side of the
  split plane.
\end{definition}

\subsection{Theory}

The query on an orthogonal range $Q$ descends the tree: if a node's cell
does not intersect $Q$, the whole subtree is pruned; if the cell is
contained in $Q$, the whole subtree is reported; otherwise the search
continues in both children. The number of nodes visited is bounded by
$O(n^{1-1/d} + k)$ in the worst case for a box query in $\mathbb{R}^d$
with $n$ points; the exponent reflects that pruning eliminates entire
subtrees only in dimensions with ``thick'' splits. The $k$-d tree is the
structure of choice when the dimension is small ($d \le 3$) and the points
are static or slowly changing, because its storage is exactly linear and
its constants are tiny.

\subsection{Pseudocode}

\begin{algorithm}[htbp]
\caption{Orthogonal Range Query in a $k$-d Tree}
\label{alg:kd_range_query}
\begin{algorithmic}[1]
\Require $k$-d tree root with cells, orthogonal query box $Q$.
\Ensure All points of the tree inside $Q$.
\Function{KDQuery}{$v, Q$}
  \If{$v$ is a leaf} \State \Return $v.\text{point}$ if it lies in $Q$
  \EndIf
  \If{$\text{cell}(v) \cap Q = \emptyset$} \State \Return $\emptyset$ \EndIf
  \If{$\text{cell}(v) \subseteq Q$} \State \Return all points of the subtree \EndIf
  \State \Return $\textsc{KDQuery}(\text{left}(v), Q) \cup \textsc{KDQuery}(\text{right}(v), Q)$
\EndFunction
\end{algorithmic}
\end{algorithm}

\subsection{Complexity}

\begin{itemize}
  \item \textbf{Storage}: $O(n)$.
  \item \textbf{Construction}: $O(n \log n)$ by median finding at each node.
  \item \textbf{Query}: $O(n^{1-1/d} + k)$ worst case for orthogonal ranges;
    the expected behavior on typical data is much better (logarithmic plus
    output).
\end{itemize}

\subsection{Usages}

\begin{itemize}
  \item \textbf{Nearest-neighbor search}: the branch-and-bound version of
    the $k$-d tree is the workhorse of Chapter~\ref{ch:proximity}.
  \item \textbf{Windowing}: reporting all points in a query rectangle, used
    in GIS and graphics.
  \item \textbf{Dynamic point sets}: reinsertion-based updates make the
    $k$-d tree usable online (Section~\ref{sec:rangesearch_dynamic}).
\end{itemize}

\subsection{References}

The $k$-d tree was introduced by Bentley \cite{bentley1975}; the analysis
of its range queries and the branch-and-bound nearest-neighbor search are
due to Friedman, Bentley, and Finkel \cite{friedman1977}; textbook
treatments appear in \cite{berg2008,preparata1985}.

\section{Range Trees}
\label{sec:rangesearch_range_trees}

\subsection{Overview}

The \textbf{range tree}\index{range tree} answers orthogonal box queries in
the plane in $O(\log^2 n + k)$ time with $O(n \log n)$ storage, and in
$d$ dimensions in $O(\log^d n + k)$ time with $O(n \log^{d-1} n)$ storage.
Unlike the $k$-d tree, its query time is near-logarithmic in $n$ and
independent of the answer count's position on the $n^{1-1/d}$ curve; it is
the textbook example of a \textbf{multilevel} (or layered) data structure.

\subsection{Definitions}

\begin{definition}[Range Tree]\label{def:range_tree}
  A \textbf{range tree} for a set $P$ of points in $\mathbb{R}^2$ is a
  balanced binary search tree $T$ on the $x$-coordinates of $P$. Each node
  $v$ stores the subset $P_v$ of points in its subtree, organized as a
  one-dimensional balanced tree (an \textbf{associate structure}) on their
  $y$-coordinates.
\end{definition}

\subsection{Theory}

A box query $[x_1,x_2] \times [y_1,y_2]$ is answered in two stages. First,
search the primary tree for the $x$-range and collect the
\textbf{canonical subset}: the $O(\log n)$ nodes whose subtrees exactly
cover the points with $x$-coordinates in $[x_1,x_2]$. Then, for each such
node $v$, query $v$'s associate $y$-tree for the interval $[y_1,y_2]$ and
report the output. Counting is the same, with each associate query
returning a count. The two-stage structure is correct because the points
with $x \in [x_1,x_2]$ are exactly the disjoint union of the canonical
nodes' subsets, and each is filtered by the $y$-query.

\subsection{Pseudocode}

\begin{algorithm}[htbp]
\caption{2D Orthogonal Range Query}
\label{alg:range_tree_query}
\begin{algorithmic}[1]
\Require Range tree $T$ on $P$, box $Q = [x_1,x_2] \times [y_1,y_2]$.
\Ensure All points of $P$ in $Q$.
\Function{RangeTreeQuery}{$T, Q$}
  \State collect the $O(\log n)$ canonical nodes $v_1,\dots,v_m$ covering $[x_1,x_2]$
  \For{each canonical node $v_i$}
    \State report the points of $v_i$'s $y$-tree with $y \in [y_1,y_2]$
  \EndFor
\EndFunction
\end{algorithmic}
\end{algorithm}

\subsection{Complexity}

\begin{itemize}
  \item \textbf{Storage}: $O(n \log n)$; each point appears in one
    associate tree per level, of which there are $\log n$.
  \item \textbf{Construction}: $O(n \log n)$ by building the associate
    trees bottom-up.
  \item \textbf{Query}: $O(\log^2 n + k)$ in 2D; the canonical search costs
    $O(\log n)$ and each of the $O(\log n)$ associate queries costs
    $O(\log n)$.
\end{itemize}

\subsection{Usages}

\begin{itemize}
  \item \textbf{Orthogonal windowing}: the canonical-subset mechanism is
    the standard answer for exact-logarithmic queries on static data.
  \item \textbf{Dual problems}: counting points in a query box is the
    direct analogue of the database ``box select''.
  \item \textbf{Reduction to 1D}: any problem solvable on a sorted set in
    $O(\log n)$ transfers to $O(\log^2 n)$ in the plane via the range tree,
    which is how the range tree is used as a black box inside other
    algorithms.
\end{itemize}

\subsection{References}

The range tree (as the ``layered range tree'' using multidimensional
divide-and-conquer) is due to Bentley \cite{bentley1980multidim}; the
higher-dimensional analysis and the fractional-cascading speedup are
developed in \cite{berg2008} and \cite{willard1985}.

\section{Multilevel Data Structures}
\label{sec:rangesearch_multilevel}

A \textbf{multilevel data structure}\index{multilevel data structure}
answers a query by filtering through a hierarchy of one-dimensional
structures, one level per dimension. The range tree is the canonical
example: the primary level filters on $x$, and each canonical node contains
a secondary structure that filters on $y$. The construction recurses: a
$d$-dimensional range tree has a primary tree on the first coordinate and,
at each node, a $(d-1)$-dimensional range tree on the remaining
coordinates, with the base case $d=1$ being the balanced tree of
Section~\ref{sec:rangesearch_1d}.

The general design principle is: if a query on a set can be answered in
$Q(n)$ time and the associate structures can be built in $S(n)$ storage per
node, then the two-level structure answers in $O(\log n)$ (canonical
search) times $Q(n)$ and stores $O(n)$ times the per-node storage. Applied
recursively, this yields the $d$-dimensional bounds of
Section~\ref{sec:rangesearch_higher_dimensional}. The same principle
underlies many geometric data structures beyond range searching; it is
sometimes called \textbf{multidimensional divide-and-conquer}
\cite{bentley1980multidim}.

Two implementation notes matter in practice. First, the associate structures
at a node contain the points of the subtree, so a point participates in one
structure per level, giving the $\log n$ blow-up in storage. Second, when
the query filter is a contiguous interval in the associated coordinate, the
associate structure can be a sorted array and the query becomes two binary
searches --- this is the form used when fractional cascading
(Section~\ref{sec:rangesearch_fractional_cascading}) is applied.

\index{multilevel structure}\index{multidimensional divide-and-conquer}

\section{Fractional Cascading}
\label{sec:rangesearch_fractional_cascading}

\subsection{Overview}

A 2D range-tree query spends $O(\log n)$ time finding the canonical nodes
and then performs $O(\log n)$ binary searches, one per canonical node, for a
total of $O(\log^2 n)$. Each binary search looks for the same $y$-value in a
related sorted list, and the repeated searches are wasteful: the list at a
node and the list at its children are highly correlated. \textbf{Fractional
cascading}\index{fractional cascading} removes the second $\log n$ factor
by passing the search position from each list to the next, reducing the
query time to $O(\log n + k)$.

\subsection{Definitions}

\begin{definition}[Fractional Cascading]\label{def:rangesearch_cascading}
  \textbf{Fractional cascading} is a technique for sequences of sorted
  lists $L_1, L_2, \dots$ where each list's successor can be searched
  faster by \emph{cascading}: every element of $L_{i+1}$ stores a pointer
  to the closest elements in $L_i$, so that a query value's position in
  $L_i$ determines its position in $L_{i+1}$ in $O(1)$ time rather than
  $O(\log n)$.
\end{definition}

\subsection{Theory}

Each $y$-list in the range tree is replaced by a \textbf{cascading array}:
for every element, a pair of pointers to the nearest elements of the same
value range in the lists of the node's children. A query that has found
its position in the root's $y$-list then walks down the tree, spending
$O(1)$ per level to derive the position in the child's list. Searching all
$O(\log n)$ canonical $y$-lists therefore costs $O(\log n)$ total instead
of $O(\log^2 n)$. Because the lists are now linked, the storage remains
$O(n \log n)$ but the constant grows by a small factor (each element keeps
two extra pointers).

\subsection{Complexity}

\begin{itemize}
  \item \textbf{Query}: $O(\log n + k)$ for 2D orthogonal boxes.
  \item \textbf{Storage}: $O(n \log n)$.
  \item \textbf{Construction}: $O(n \log n)$.
\end{itemize}

\subsection{Usages}

\begin{itemize}
  \item \textbf{Optimal orthogonal searching}: together with careful
    canonical decomposition, fractional cascading gives the optimal
    $O(\log n + k)$ 2D structure of Willard \cite{willard1985}.
  \item \textbf{Any repeated-scan hierarchy}: the technique applies
    wherever the same value is searched in related sorted lists, which
    occurs in several other chapters (e.g., computing lower envelopes).
\end{itemize}

\subsection{References}

Fractional cascading was introduced by Chazelle and Guibas \cite{chazelle1986cascading};
its application to range trees is the textbook treatment of \cite{berg2008}
and the optimal-structure result of \cite{willard1985}.

\section{Priority Search Trees}
\label{sec:rangesearch_priority_search}

\subsection{Overview}

A \textbf{priority search tree}\index{priority search tree} \cite{mccreight1985}
answers \textbf{3-sided} orthogonal queries --- ranges of the form
$[x_1,x_2] \times (-\infty, y]$ --- in $O(\log n + k)$ time with only
$O(n)$ storage. The improvement over the range tree's $O(n \log n)$
storage comes from combining the $x$-order (tree structure) and the
$y$-order (heap order) in a single binary tree, a ``heap on $y$ over a
search tree on $x$''.

\subsection{Definitions}

\begin{definition}[Priority Search Tree]\label{def:priority_search_tree}
  A \textbf{priority search tree} for a set $P$ of $n$ points is a binary
  tree where the root holds the point of minimal $y$, the remaining points
  are split by $x$-median into the left and right subtrees, and each
  subtree is itself a priority search tree. Every node stores the
  $x$-range (min and max $x$) of its subtree.
\end{definition}

\subsection{Theory}

A 3-sided query $[x_1,x_2] \times (-\infty, y]$ is answered by a pruning
search: a node is visited only if its $x$-range intersects the query
interval. When a visited node's point has $y$-coordinate at most $y$, the
point is reported (the heap order guarantees that if the root of a subtree
fails the $y$-test, no point of that subtree passes it). The number of
visited nodes is $O(\log n + k)$, giving the near-logarithmic bound with
linear storage. An axis-aligned rectangle query can be decomposed into two
3-sided queries, so priority search trees also answer full orthogonal
range queries, though with $O(k)$ reporting and $O(\log n)$ counting
structure trade-offs different from the range tree.

\subsection{Complexity}

\begin{itemize}
  \item \textbf{Storage}: $O(n)$.
  \item \textbf{Query (3-sided)}: $O(\log n + k)$.
  \item \textbf{Insertion/deletion}: $O(\log n)$.
\end{itemize}

\subsection{Usages}

\begin{itemize}
  \item \textbf{Dynamic 3-sided queries}: the heap order supports cheap
    insertions, making the structure attractive when points arrive online.
  \item \textbf{Component of box structures}: 3-sided queries are a
    standard primitive into which general orthogonal queries are
    decomposed.
\end{itemize}

\subsection{References}

The structure is due to McCreight \cite{mccreight1985}; its relation to
the range tree and its use in dynamic structures are covered in
\cite{berg2008} and \cite{overmars1983}.

\section{Reporting versus Counting}
\label{sec:rangesearch_reporting_counting}

Range queries come in two flavors whose complexity differs in an important
way. \textbf{Reporting} queries must list the $k$ points in the range, so
any structure costs at least $\Omega(k)$ per query, and a structure with
$O(\log n)$ search time can report in $O(\log n + k)$. \textbf{Counting}
queries return only the number $|P \cap Q|$, so their time is independent
of $k$; the challenge is that counting often requires storing more
information per canonical subset (a subtree size, or a precomputed count)
while spending as little time as possible inside the structure.

The two versions have different optimality landscapes:
\begin{itemize}
  \item for counting, a range tree stores the size of each $y$-tree, so a
    box count costs $O(\log^d n)$ with $O(n \log^{d-1} n)$ storage;
  \item for reporting, the same tree costs $O(\log^d n + k)$, where $k$
    includes the output; priority search trees and $k$-d trees trade
    query time for output time differently;
  \item lower bounds differ as well: counting is easier than reporting in
    the sense that the output term is absent, and decision-tree arguments
    give $O(n \log n)$-ish bounds for the static 2D case, while reporting
    has the trivial $\Omega(k)$ component.
\end{itemize}
In practice most applications need reporting, and the $k$ term dominates;
counting appears in statistics (counts of points in half-spaces for
classification) and in the simplex setting of Section~\ref{sec:rangesearch_simplex},
where the near-linear-space lower bound $\Omega(n^{1-1/d})$ applies to
counting and reporting alike.

\index{counting query}\index{reporting query}

\section{Higher-Dimensional Range Searching}
\label{sec:rangesearch_higher_dimensional}

The range-tree construction generalizes directly to $\mathbb{R}^d$: a
$d$-dimensional range tree has a primary tree on the first coordinate and
a $(d-1)$-dimensional range tree as each node's associate structure
(Section~\ref{sec:rangesearch_multilevel}). The query cost and storage
follow by induction:
\begin{itemize}
  \item \textbf{Range tree}: $O(\log^d n + k)$ query time and
    $O(n \log^{d-1} n)$ storage in $d$ dimensions.
  \item \textbf{$k$-d tree}: $O(n^{1-1/d} + k)$ query time with $O(n)$
    storage in $d$ dimensions.
  \item \textbf{Priority search tree}: $O(n \log n)$-free --- it answers
    3-sided queries in $\mathbb{R}^2$ in $O(\log n + k)$ with $O(n)$
    storage; in higher dimensions its role is taken by more specialized
    structures.
\end{itemize}
The range tree's storage grows polynomially in the dimension (the
$n\log^{d-1} n$ factor), so for $d$ fixed it is the near-logarithmic
structure of choice; for variable dimension the polynomial storage becomes
prohibitive. The $k$-d tree, with its strictly linear storage, is the
practical default when $d$ is moderate (say $d \le 10$), because the
$n^{1-1/d}$ query factor is small until $d$ approaches the data size. The
extreme range --- where every structure degenerates toward linear scans ---
is the curse of dimensionality of Section~\ref{sec:rangesearch_curse}.
Simplex range searching in higher dimensions, with its near-optimal
$O(n^{1-1/d}(\log n)^{O(1)})$ bounds, is treated in
Section~\ref{sec:rangesearch_simplex}.

\section{Dynamic Range Searching}
\label{sec:rangesearch_dynamic}

When points are inserted and deleted, a range-searching structure must
maintain its invariants online. The classical framework is the
transform-and-conquer method of Overmars \cite{overmars1983}, which
converts any static structure into a dynamic one at the cost of a
logarithmic factor per operation, provided the static structure can be
rebuilt and queried efficiently. The main strategies are:

\begin{description}
  \item[Binary-search-tree dynamics] A dynamic range tree keeps the
    primary tree balanced with rotations; each rotation changes only
    $O(1)$ associate structures locally, so insertions and deletions cost
    $O(\log^2 n)$ in 2D (an $O(\log n)$ factor for the tree update times
    an $O(\log n)$ factor for the associate-structure updates).
  \item[Rebuilding] Overmars' method maintains several static structures
    of geometrically growing sizes (partial rebuilding): a new point is
    inserted into the smallest structure, and structures are merged when
    they exceed a size threshold. Querying all levels costs $O(\log n)$
    times the static query cost.
  \item[$k$-d tree updates] Insertions into a $k$-d tree are handled by
    attaching a point at the appropriate leaf; deletions mark points and
    periodically rebuild the subtree. The amortized update cost is
    $O(\log n)$ and queries remain $O(n^{1-1/d} + k)$.
\end{description}

The fundamental limitation is that updates interfere with fractional
cascading and heap orders (Section~\ref{sec:rangesearch_priority_search}),
so the fully dynamic optimal 2D structure is considerably more complex
than the static one; in practice, the field's recommendation is to use a
dynamic range tree for moderate update rates and a $k$-d tree when updates
are frequent and the dimension is low.

\index{dynamic data structure}\index{partial rebuilding}

\section{Curse of Dimensionality}
\label{sec:rangesearch_curse}

The \textbf{curse of dimensionality}\index{curse of dimensionality} is the
phenomenon that all data structures for orthogonal range searching
deteriorate as the dimension $d$ grows. The $k$-d tree's query factor
$n^{1-1/d}$ approaches $n$ as $d \to \infty$, so for large $d$ the query
degenerates to a linear scan; the range tree's storage
$n\log^{d-1} n$ becomes astronomical in $d$; and, most fundamentally, the
lower bounds of Section~\ref{sec:rangesearch_simplex} show that any
structure with near-linear storage must spend $\Omega(n^{1-1/d})$ time per
query in the worst case. For $d = \log n$, the exponents meet: every
near-linear-space structure answers queries in time at least $n^{1-1/\log
n} = n^{1 - o(1)}$, i.e., no better than scanning.

The cause is combinatorial: the number of distinct orthogonal ranges that
a point set can define grows exponentially in $d$, so a structure that
wants to answer all of them with few canonical subsets cannot exist. Two
practical responses are standard. First, for moderate $d$ ($d \le 20$),
indexing schemes such as the $k$-d tree and the space-filling-curve index
of Chapter~\ref{ch:space_filling_curves} still perform well on real data,
whose intrinsic dimension is far below $d$; the analysis in
Chapter~\ref{ch:high_dimensional} explains when this holds. Second, when
$d$ is genuinely large, exact range searching is abandoned in favor of
approximate structures --- locality-sensitive hashing, approximate nearest
neighbors, and sketches --- whose guarantees are probabilistic rather than
worst-case (Chapter~\ref{ch:high_dimensional}).

\section{Simplex Range Searching}
\label{sec:rangesearch_simplex}

\subsection{Overview}
Range searching asks for the set of points of a fixed point set $P$ that lie
inside a query region $Q$. When $Q$ is a simplex (a half-plane, triangle, or
half-space), the problem is called \textbf{simplex range searching}. It is one
of the most fundamental query problems in computational geometry, with
applications in databases, GIS, and computer graphics. The standard solutions
use \textbf{partition trees}, \textbf{cuttings}, and \textbf{$\varepsilon$-nets}
to answer counting queries in roughly $O(n^{1-1/d})$ time with linear storage,
which is optimal in the worst case for near-linear space.

\subsection{Definitions}

\begin{definition}[Simplex Range Counting]\label{def:simplex_range_counting}
Given a set $P$ of $n$ points in $\mathbb{R}^d$ and a query simplex $\Delta$,
the \emph{simplex range counting} problem asks for $|P \cap \Delta|$. The
\emph{reporting} variant asks for all points in $P \cap \Delta$.
\end{definition}

\begin{definition}[$\varepsilon$-Net]\label{def:eps_net}
Let $(X, \mathcal{R})$ be a range space of points and subsets (ranges) with VC
dimension $d$. A subset $N \subseteq X$ is an \emph{$\varepsilon$-net} if every
range that contains at least $\varepsilon |X|$ points also contains a point of
$N$.
\end{definition}

\subsection{Theory}
A random sample of size $O(\frac{d}{\varepsilon}\log\frac{1}{\varepsilon})$ is
an $\varepsilon$-net with high probability~\cite{hausslerwelzl1987}. Cuttings
generalize this idea: a \emph{$(1/r)$-cutting} of an arrangement of $n$ lines
subdivides the plane into $O(r^2)$ cells, each crossed by at most $n/r$ lines.
Cuttings can be built deterministically in $O(nr)$ time~\cite{chazellefriedman1990},
and hierarchical cuttings yield partition trees.

\begin{theorem}[Simplex Range Searching Bounds]\label{thm:simplex_range_searching}
For a set of $n$ points in $\mathbb{R}^d$, a partition tree of linear size can
answer a simplex counting query in $O(n^{1-1/d}\,(\log n)^{O(1)})$ time and a
reporting query in $O(n^{1-1/d}\,(\log n)^{O(1)} + k)$ time, where $k$ is the
number of reported points~\cite{matousek1992,chazellefriedman1990}.
This is optimal for near-linear space: any data structure using $O(n
\operatorname{polylog} n)$ space requires $\Omega(n^{1-1/d})$ query time in
the worst case.
\end{theorem}

\subsection{Pseudocode}

\begin{algorithm}[ht]
\label{alg:partition_tree}
\caption{Simplex Range Counting with a Partition Tree}
\begin{algorithmic}[1]
\Function{CountSimplex}{node, $\Delta$}
  \State \Comment{$node$ is a cell with an associated canonical subset}
  \If{$node.\text{cell} \cap \Delta = \emptyset$}
    \State \Return $0$
  \EndIf
  \If{$node.\text{cell} \subseteq \Delta$}
    \State \Return $node.\text{count}$
  \EndIf
  \State \Return $\sum_{child} \textsc{CountSimplex}(child, \Delta)$
\EndFunction
\end{algorithmic}
\end{algorithm}

\subsection{Complexity}
\begin{itemize}
    \item \textbf{Storage}: $O(n)$ for a partition tree; cutting-based
      structures use $O(n \operatorname{polylog} n)$.
    \item \textbf{Counting Query}: $O(n^{1-1/d}\,(\log n)^{O(1)})$ time.
    \item \textbf{Reporting Query}: add $O(k)$ output time.
    \item \textbf{Construction}: $O(n \log n)$ with randomized incremental
      methods, or $O(nr)$ for a single cutting.
\end{itemize}

\subsection{Usages}
\begin{itemize}
    \item \textbf{Database and GIS} polygon windowing queries.
    \item \textbf{Geometric statistics}: counting points in a query half-space
      for classification and density estimation.
    \item \textbf{Levels and duality}: by duality, half-plane range counting
      is equivalent to counting vertices of an arrangement below the query
      line, connecting to levels (Section~\ref{sec:arrangement_levels}).
\end{itemize}

\subsection{Testing Methods and Edge Cases}
\begin{itemize}
    \item \textbf{Degenerate queries}: simplex boundaries through points,
      empty and full queries.
    \item \textbf{Consistency}: compare counting against brute force for small
      $n$; compare reporting to counting by checking the output size.
\end{itemize}

\subsection{References}
\begin{itemize}
    \item Haussler, D., \& Welzl, E. (1987). ``$\varepsilon$-Nets and Simplex
      Range Queries''. Discrete \& Computational Geometry.
    \item Matou\v{s}ek, J. (1992). ``Reporting Points in Halfspaces''
      Computational Geometry: Theory and Applications.
    \item Chazelle, B., \& Friedman, J. (1990). ``A Deterministic View of
      Random Sampling and Its Use in Geometry''. Combinatorica.
    \item Matou\v{s}ek, J. (1999). ``Geometric Discrepancy: An Illustrated
      Guide''. Springer.
\end{itemize}

\section{Exercises}
\label{sec:rangesearch_exercises}

\begin{enumerate}
  \item Show that a 1D range tree (Section~\ref{sec:rangesearch_1d}) is
    just a sorted array, and that counting needs only the binary-search
    indices.
  \item Prove that the $k$-d tree query visits $O(n^{1-1/d} + k)$ nodes in
    the worst case for a box query: argue that a query box whose side is
    unbounded prunes half the subtrees per level, while a thin box cuts
    many subtrees.
  \item Verify the range-tree bounds: each point is stored in exactly one
    associate tree per level, so storage is $O(n \log n)$, and the query
    visits $O(\log n)$ canonical nodes.
  \item Explain why fractional cascading reduces the query time by a
    factor of $\log n$ and why the storage stays $O(n \log n)$.
  \item Design a priority search tree that answers the query
    $[x_1, x_2] \times (-\infty, y]$ and prove the visited-node bound
    $O(\log n + k)$.
  \item Compare the asymptotic costs of range trees and $k$-d trees for
    orthogonal reporting in $d = 2, 3, 10$, and state which you would
    choose for each.
  \item Argue from the lower bound of Section~\ref{sec:rangesearch_simplex}
    that no near-linear-space structure can answer counting queries
    faster than $n^{1-1/d}$ in the worst case, and interpret the bound for
    $d = \log n$.
  \item Implement a brute-force counter and verify a range tree on random
    point sets, checking that reporting agrees with brute force and that
    counting equals the reported cardinality.
\end{enumerate}

\section{Project: Spatial Query Engine}
\label{sec:rangesearch_project}

Build a spatial query engine over a point dataset:
\begin{enumerate}
  \item \textbf{Index}: implement a 2D range tree (primary tree on $x$,
    associate sorted arrays on $y$) and a $k$-d tree.
  \item \textbf{Queries}: support box reporting and box counting, and a
    point-in-polygon query by bounding-box filtering followed by an exact
    test.
  \item \textbf{Verification}: compare both indexes against a brute-force
    scan on random and adversarial point sets, checking output equality and
    reporting counts.
  \item \textbf{Measurement}: plot query time and storage against $n$ and
    against the query box size; verify the predicted $O(\log^2 n + k)$ and
    $O(n^{1-1/2} + k)$ behaviors.
  \item \textbf{Extras}: add the 3-sided query with a priority search
    tree, and a nearest-neighbor query reusing the $k$-d tree
    (Chapter~\ref{ch:proximity}).
  \item \textbf{Visualization}: draw the points, the query boxes, and the
    reported points; optionally animate the recursive subdivision of the
    $k$-d tree.
\end{enumerate}
